We now assess our method's precision-recall performance at mortality risk predictions on two large open-access datasets of critical care EHR: MIMIC-III (Medical Information Mart for Intensive Care) and eICU.
Our goal is to predict in-hospital mortality after observing vitals and labs over the entire sequence of a patient's stay, up to just before discharge or death.
We enforce a minimum precision value of $\alpha = 0.9$ on the training set.
We found $\alpha=0.9$ challenging to meet on heldout data even when satisfied at training, so to select the best runs we enforce a minimum validation precision $\alpha' = 0.8$.

\textbf{MIMIC-III data}. The MIMIC-III dataset \citep{johnsonMIMICIIIFreelyAccessible2016} contains deidentified open-access EHR data from a Boston hospital from 2002-2012. We extract 2 demographics, 10 vitals, and 94 lab measurements discretized to hourly bins for 34472 patient stays using the MIMIC-Extract pipeline~\citep{wangMIMICExtractDataExtraction2020}.  
Our train/valid./test splits have 20682/6895/6895 patient-stays, with $\sim 9.5\%$ patient stays resulting in death. 
Each patient-stay's multivariate time-series is transformed into a feature vector as follows. For each raw feature, we apply 7 summarization functions (missing indicator; time since last non-missing; plus min, max, median, slope, and variance of non-missing values) to 4 possible time windows 
(0-100\%, 50\%-100\%, last 16 hr, and last 24 hr). 
The resulting feature vector has size 106x7x4=2968.

\textbf{eICU data.}
The eICU Collaborative Research Database \citep{pollard2018eicu} contains data from 59 critical care units throughout the U.S. We extract 3 demographics, 8 vitals, and 6 lab measurements discretized to hourly bins using eICU Extract~\citep{eICUExtract} for 72670 patient stays.
Our train/valid/test splits have 43642/14509/14518 patient-stays, with $\sim 8.2\%$ resulting in death.
We featurize each patient-stay the same way as MIMIC, yielding feature vectors of size 17x7x4=417.

\textbf{Models.}
On both datasets, we tried two models: logistic regression (LR) and a multi-layer perceptron (MLP) with 1 hidden layer (32 hidden units; RELU activation).
We search 25 random seeds, 2 initialization scale factors, 2 $\lambda$ values, and 5 weight decays.

\textbf{Mortality Prediction Results.}
The recall performance of each method on both datasets can be found in Table~\ref{table:PR_perseq_mimic_and_eicu}.
Across 20 independent randomly-chosen train-test partitions of both datasets, we find that all methods can satisfy the desired precision constraint (all methods exceed 0.9 on training and 0.8 on validation). However, our proposed sigmoid bound method consistently achieves better test-set recall values at the selected operating threshold (chosen on validation data).
Compared to the BCE baseline, using a linear model we see our method boost the test-set recall (median across partitions) from 0.540 to 0.684 on MIMIC and from 0.087 to 0.195 on eICU.
As expected, gains from the MLP are smaller (as a predictor gets more flexible, it can get closer to the optimal non-linear boundary).
However, the gains are still noticeable: recall improves from 0.689 to 0.717 on MIMIC and from 0.280 to 0.298 on eICU.
\citet{ebanScalableLearningNonDecomposable2017}'s hinge bound method is \emph{worse} than the BCE baseline in terms of test-set recall on three of the four dataset-model combinations tested.

Across Table~\ref{table:PR_perseq_mimic_and_eicu}, recall values in on eICU are noticeably lower than MIMIC (BCE achieves test recall of 0.5-0.7 on MIMIC but only 0.08-0.28 for eICU). This occurs despite the desired minimum precision value remaining the same ($\alpha = 0.9$) and being satisfied in all cases.
The drop in recall performance likely occurs because we have more available informative clinical measurements for MIMIC than eICU (106 vs. 17). 
The eICU task is also slightly more difficult: only 8\% of eICU sequences have positive labels compared to 9.5\% for MIMIC. 
In this challenging setting, enforcing high precision means some modest tradeoffs with recall. We emphasize that without enforcing sufficiently high precision, alarm fatigue could render the whole system ineffective.

%The loose hinge bound has especially poor recall in this regime. In our imagined use case, all patients still receive routine hospital care, even those missed by our alerts. Avoiding alarm fatigue is critical}


To better understand the gains of our method in terms of common evaluation criteria, we present ROC and PR curves for all methods in the appendix. Our method shows consistent gains in both curves.

% We do stress that measuring area under the ROC curve is inadequate for measuring the quality of the alerts produced by a specific operating threshold of an early warning system~\citep{romero-brufauWhyCstatisticNot2015}, which is why we favor the precision and recall metrics in Table~\ref{table:PR_perseq_mimic_and_eicu}.
